\documentclass[12pt,a4paper]{article}

\usepackage[utf8]{inputenc}
\usepackage[T1]{fontenc}
\usepackage[ngerman]{babel}
\usepackage{geometry}
\geometry{margin=2.5cm}
\usepackage{setspace}
\onehalfspacing
\usepackage{graphicx}
\usepackage{longtable}
\usepackage{hyperref}
\usepackage{enumitem}
\usepackage{titlesec}
\usepackage{xcolor}
\usepackage{fancyhdr}

\pagestyle{fancy}
\fancyhf{}
\rhead{Projektbericht -- Serious Game Plattform}
\lhead{\leftmark}
\cfoot{\thepage}

\title{Projektbericht\\[0.4em]\large Weiterentwicklung einer Serious-Game-Plattform im Gesundheitskontext}
\author{Max Mustermann\\Projektzeitraum: vier Monate}
\date{\today}

\begin{document}

\maketitle
\thispagestyle{empty}
\newpage

\tableofcontents
\newpage

\section{Einleitung}
Digitale Anwendungen gewinnen im Gesundheitskontext zunehmend an Bedeutung. Neben klassischen Web- und Mobile-Anwendungen entstehen immer mehr interaktive Systeme, die Training, Motivation und Verlaufskontrolle miteinander verbinden. In diesem Umfeld nehmen \textbf{Serious Games} eine besondere Rolle ein: Sie verknüpfen spielerische Elemente mit funktionalen Zielen wie Rehabilitation, Aktivierung, Therapiebegleitung oder strukturierter Datenerfassung.

Das vorliegende Projekt befasst sich mit der Weiterentwicklung einer Serious-Game-Plattform, die mehrere technische Bausteine in einem gemeinsamen System zusammenführt. Dazu gehören ein Backend zur Daten- und Benutzerverwaltung, eine Weboberfläche für administrative und therapeutische Prozesse, ein Android-Launcher zur Bereitstellung mobiler Spiele sowie eine Test-App zur Überprüfung der Session- und Score-Übertragung.

Ziel dieses Berichts ist es, das Projekt inhaltlich als \textbf{Projektbericht} zu beschreiben. Im Vordergrund stehen daher nicht organisatorische Rahmenbedingungen eines Praktikums, sondern die Fragen, was das Projekt fachlich und technisch aussagt, welchen aktuellen Stand die Plattform erreicht hat und welche konkreten Beiträge innerhalb eines Zeitraums von vier Monaten umgesetzt wurden.

\section{Projektübersicht}
\subsection{Was ist das Projekt?}
Die Serious-Game-Plattform ist ein verteiltes Softwaresystem zur Verwaltung, Auslieferung und Auswertung digitaler Therapie- und Trainingsspiele. Der zentrale Gedanke besteht darin, verschiedene Rollen und Nutzungssituationen innerhalb einer Plattform zusammenzubringen:

\begin{itemize}[leftmargin=1.5em]
  \item \textbf{Administration:} Spiele, Metadaten und Zuordnungen werden zentral verwaltet.
  \item \textbf{Therapeutischer Kontext:} Relevante Spiele können bestimmten Patientinnen und Patienten zugewiesen und fachlich begleitet werden.
  \item \textbf{Patientenseite:} Zugewiesene Spiele werden auf dem Zielgerät bereitgestellt und gestartet.
  \item \textbf{Auswertung:} Spielsessions, Scores und weitere Metadaten werden zurück an die Plattform übermittelt und sichtbar gemacht.
\end{itemize}

Die Plattform ist damit nicht nur eine Sammlung einzelner Apps, sondern eine durchgängige Infrastruktur für den Lebenszyklus eines Serious Games: von der technischen Einbindung über die Zuordnung bis zur Nutzung und Ergebnisdarstellung.

\subsection{Architektonische Einordnung}
Die betrachtete Lösung besteht aus mehreren Komponenten, die zusammenarbeiten:

\begin{itemize}[leftmargin=1.5em]
  \item \textbf{Backend:} Verantwortlich für Persistenz, Rollenlogik, Upload-Prozesse, Zuweisungen und Session-Daten.
  \item \textbf{Webanwendung:} Oberfläche für Admin- und Verwaltungsprozesse, insbesondere für das Hinzufügen und Verwalten von Spielen.
  \item \textbf{Android-Launcher:} Mobile Anwendung zur Anzeige und zum Start zugewiesener Android-Spiele.
  \item \textbf{Test-App:} Technische Hilfsanwendung zur Überprüfung von Session-Übertragung, Score-Rückgabe und Integrationspfaden.
\end{itemize}

Gerade die Kombination dieser Bausteine macht das Projekt relevant. Die Plattform verbindet fachliche Anforderungen aus dem Gesundheitsumfeld mit Full-Stack-Entwicklung, mobilen Bereitstellungsprozessen und einer nachvollziehbaren Datenerfassung.

\section{Aussage des Projekts}
Inhaltlich zeigt das Projekt, dass Serious Games im Gesundheitskontext nur dann sinnvoll skalierbar sind, wenn sie nicht isoliert als Einzelanwendungen betrachtet werden. Entscheidend ist eine Plattformlogik, die mehrere Anforderungen gleichzeitig erfüllt:

\begin{itemize}[leftmargin=1.5em]
  \item Spiele müssen technisch sauber eingebunden und verwaltet werden können.
  \item Die Bereitstellung auf mobilen Geräten darf nicht von fehleranfälligen manuellen Einzelschritten abhängen.
  \item Ergebnisse aus den Spielen müssen strukturiert erfasst und den relevanten Rollen wieder sichtbar gemacht werden.
  \item Die Gesamtarchitektur muss erweiterbar bleiben, damit weitere Spiele und weitere Gerätetypen angebunden werden können.
\end{itemize}

Die zentrale Aussage des Projekts lautet daher: \textbf{Eine Serious-Game-Plattform entfaltet ihren praktischen Nutzen erst durch die verlässliche Verbindung von Verwaltung, Auslieferung, Nutzung und Auswertung.} Genau an dieser Schnittstelle setzt die Weiterentwicklung an.

\section{Aktueller Stand des Projekts}
Zum aktuellen Stand liegt eine Plattform vor, deren Kernbausteine funktional miteinander verbunden sind. Besonders relevant ist dabei die inzwischen deutlich bessere Verzahnung zwischen Upload-Prozess, mobiler Bereitstellung und Anzeige von Spielergebnissen.

Der derzeitige Stand lässt sich wie folgt zusammenfassen:

\begin{itemize}[leftmargin=1.5em]
  \item Es existiert eine zentrale Struktur zur Verwaltung von Serious Games und zugehörigen Metadaten.
  \item Mobile Spiele können in den Plattformkontext integriert werden, insbesondere über APK-basierte Prozesse.
  \item Der Android-Launcher bildet die patientennahe Ausführungsebene für mobile Spiele.
  \item Sessions und Scores können an die Plattform zurückgegeben und in der Oberfläche sichtbar gemacht werden.
  \item Für Integrations- und Funktionsprüfungen steht eine Test-App als technisches Hilfsmittel zur Verfügung.
\end{itemize}

Aus Projektsicht ist die Plattform damit nicht mehr nur im Stadium einzelner Prototypen, sondern besitzt bereits einen zusammenhängenden End-to-End-Charakter. Gleichzeitig bleibt sie in mehreren Bereichen ausbaufähig, etwa bei der weiteren Automatisierung, einer noch robusteren Metadatenanalyse und umfassenderen Testabdeckung.

\section{Mein Beitrag innerhalb von vier Monaten}
\subsection{Monat 1: Einarbeitung und Systemverständnis}
Im ersten Monat lag der Schwerpunkt auf dem Verständnis der bestehenden Plattform. Dafür wurden die vorhandenen Komponenten, ihre Verantwortlichkeiten und ihre Schnittstellen zueinander analysiert. Ziel war es, die Plattform nicht nur auf Dateiebene zu verstehen, sondern den tatsächlichen fachlichen Ablauf von Spielverwaltung, Zuweisung und Nutzung nachzuvollziehen.

Schwerpunkte dieser Phase:
\begin{itemize}[leftmargin=1.5em]
  \item Analyse der Architektur aus Backend, Webanwendung, Android-Komponenten und Testpfaden.
  \item Nachvollzug des Datenflusses zwischen Upload, Assignment, Session-Übertragung und Anzeige.
  \item Identifikation manueller und potenziell fehleranfälliger Schritte im mobilen Bereitstellungsprozess.
  \item Aufbau einer belastbaren Grundlage für spätere Änderungen über mehrere Systemteile hinweg.
\end{itemize}

\subsection{Monat 2: Backend- und Integrationslogik}
Im zweiten Monat stand die technische Weiterentwicklung der Integrationslogik im Vordergrund. Dabei ging es vor allem um Abläufe, die mobile Spiele in den Plattformkontext einbinden und deren Daten später wieder verarbeitbar machen.

Wesentliche Tätigkeiten:
\begin{itemize}[leftmargin=1.5em]
  \item Überprüfung und Stabilisierung von API-nahen Abläufen im Serious-Game-Bereich.
  \item Vereinheitlichung von Validierungs- und Fehlerpfaden für relevante Artefakte.
  \item Verbesserung der Nachvollziehbarkeit im Umgang mit Spiel-Metadaten.
  \item Vorbereitung robusterer Integrationspfade für APK-basierte Spiele.
\end{itemize}

In dieser Phase war besonders wichtig, technische Änderungen so anzulegen, dass sie nicht nur lokal funktionieren, sondern in den bestehenden Plattformfluss passen.

\subsection{Monat 3: Mobile Bereitstellung und Android-Launcher-Kontext}
Im dritten Monat verlagerte sich der Fokus auf die mobile Seite des Projekts. Ziel war es, den Weg von einem eingebundenen mobilen Spiel bis zur sichtbaren Verfügbarkeit im Android-Launcher konsistenter und praxisnäher zu gestalten.

Dabei standen folgende Punkte im Mittelpunkt:
\begin{itemize}[leftmargin=1.5em]
  \item Verbesserung des Upload- und Zuordnungsflusses für mobile Spielpakete.
  \item Berücksichtigung von APK-spezifischen Eigenschaften innerhalb des Plattformmodells.
  \item Reduktion manueller Zwischenschritte bei der Bereitstellung im Launcher-Kontext.
  \item Bessere Grundlage für eine direkte Nutzung neuer mobiler Spiele auf dem Zielgerät.
\end{itemize}

Diese Phase war für den Projektcharakter besonders wichtig, weil hier sichtbar wurde, wie stark technische Infrastruktur und tatsächliche Nutzbarkeit zusammenhängen.

\subsection{Monat 4: Sichtbarkeit von Sessions, Scores und Testpfaden}
Im vierten Monat lag der Schwerpunkt auf der Frage, wie Ergebnisse aus Serious Games nicht nur technisch empfangen, sondern auch sinnvoll sichtbar gemacht werden können. Gleichzeitig wurden Testpfade genutzt, um die Integrationskette zwischen Spiel, Launcher, Test-App und Plattform besser nachvollziehen zu können.

Umgesetzt beziehungsweise vertieft wurden dabei:
\begin{itemize}[leftmargin=1.5em]
  \item Sichtbarmachung von Session-Daten in der Benutzeroberfläche.
  \item Einbindung beziehungsweise Nutzung von Score- und Highscore-Informationen.
  \item Verbesserung der Nachvollziehbarkeit von Session-Metadaten.
  \item Nutzung einer Test-App zur Validierung des technischen End-to-End-Verhaltens.
\end{itemize}

Damit wurde die Plattform nicht nur um technische Funktionen ergänzt, sondern auch in ihrer Beobachtbarkeit verbessert. Gerade im Gesundheitskontext ist diese Transparenz wichtig, weil die fachliche Aussage eines Spiels wesentlich von den zurückgelieferten Ergebnissen abhängt.

\section{Fachliche und technische Ergebnisse}
Aus dem viermonatigen Projektzeitraum lassen sich mehrere konkrete Ergebnisse ableiten:

\begin{enumerate}[leftmargin=1.5em]
  \item \textbf{Klarerer Projektzusammenhang:} Die Plattform wurde als zusammenhängendes System aus Verwaltung, mobiler Bereitstellung und Auswertung weiter geschärft.
  \item \textbf{Verbesserte mobile Einbindung:} APK-bezogene Abläufe wurden stärker in den Plattformprozess integriert.
  \item \textbf{Bessere Sichtbarkeit von Ergebnissen:} Sessions und Scores sind für die Plattformnutzung relevanter und nachvollziehbarer geworden.
  \item \textbf{Unterstützung durch Testpfade:} Mit der Test-App kann das Zusammenspiel der Komponenten besser überprüft werden.
  \item \textbf{Stärkere Erweiterbarkeit:} Die Plattform wurde in eine Richtung weiterentwickelt, in der zusätzliche Spiele und mobile Szenarien einfacher anschließbar sind.
\end{enumerate}

\section{Herausforderungen}
\subsection*{Systemgrenzen und Schnittstellen}
Die größte technische Herausforderung bestand darin, Änderungen nicht isoliert zu betrachten. Schon kleine Anpassungen an Datenmodellen oder Integrationslogik wirken sich auf Backend, Weboberfläche, Launcher und Testpfade aus.

\subsection*{Mobile Bereitstellung}
Die Einbindung von Android-Spielen ist anspruchsvoller als ein reiner Upload-Prozess. Relevante Metadaten, Zuweisungen und die tatsächliche Sichtbarkeit im Launcher müssen zusammenpassen, damit aus einem Artefakt auch wirklich ein nutzbares Spiel wird.

\subsection*{Nachvollziehbarkeit von Ergebnissen}
Session- und Score-Daten sind nur dann fachlich wertvoll, wenn sie konsistent übertragen, gespeichert und angezeigt werden. Dadurch entsteht zusätzlicher Abstimmungsbedarf zwischen technischen Schichten.

\section{Fazit und Ausblick}
Der aktuelle Projektstand zeigt, dass die Serious-Game-Plattform eine belastbare Grundlage für die Verwaltung, Bereitstellung und Auswertung digitaler Therapie- und Trainingsspiele bietet. Besonders relevant ist die stärkere Verbindung zwischen administrativen Prozessen, mobiler Nutzung über den Android-Launcher und der Rückführung von Session-Daten in die Plattform.

Innerhalb von vier Monaten lag mein Beitrag vor allem in der Analyse und Weiterentwicklung genau dieser Integrationspunkte. Dadurch wurde die Plattform in einem Bereich gestärkt, der für ihren praktischen Nutzen entscheidend ist: dem Übergang von einem verwalteten Spielobjekt zu einem tatsächlich nutzbaren und auswertbaren Serious Game.

Für weitere Ausbaustufen sind insbesondere folgende Themen sinnvoll:
\begin{itemize}[leftmargin=1.5em]
  \item robustere und stärker automatisierte Metadatenanalyse für mobile Spielpakete,
  \item erweiterte Auswertungsansichten für Sessions und Scores,
  \item zusätzliche Integrations- und End-to-End-Tests,
  \item klarere Deployment- und Betriebsprozesse für mobile Spiele im Plattformkontext.
\end{itemize}

\section{Platzhalter für Screenshots}
Die folgenden Abschnitte dienen als Platzhalter für spätere visuelle Ergänzungen im Projektbericht.

\subsection*{Screenshot 1 -- Launcher Startansicht}
\fbox{\parbox[c][7cm][c]{\textwidth}{\centering Platzhalter: Android-Launcher Startansicht}}

\subsection*{Screenshot 2 -- Launcher Spieleübersicht}
\fbox{\parbox[c][7cm][c]{\textwidth}{\centering Platzhalter: Launcher mit zugewiesenen Spielen}}

\subsection*{Screenshot 3 -- Launcher Detail- oder Auswahlansicht}
\fbox{\parbox[c][7cm][c]{\textwidth}{\centering Platzhalter: Launcher-View mit Spielauswahl oder Detailansicht}}

\subsection*{Screenshot 4 -- Webansicht für Spielverwaltung}
\fbox{\parbox[c][7cm][c]{\textwidth}{\centering Platzhalter: Webansicht für Serious-Game-Verwaltung oder Upload}}

\subsection*{Screenshot 5 -- Session- und Score-Ansicht}
\fbox{\parbox[c][7cm][c]{\textwidth}{\centering Platzhalter: Anzeige von Sessions, Scores oder Highscores}}

\subsection*{Screenshot 6 -- Test-App Startansicht}
\fbox{\parbox[c][7cm][c]{\textwidth}{\centering Platzhalter: Test-App Startansicht}}

\subsection*{Screenshot 7 -- Test-App Session-Übertragung}
\fbox{\parbox[c][7cm][c]{\textwidth}{\centering Platzhalter: Test-App bei Session- oder Score-Submission}}

\newpage
\appendix
\section{Anhang: README-Hinweis zum Android-Launcher}
Der Android-Launcher ist innerhalb der Plattform die mobile Ausführungsebene für zugewiesene Spiele. Aus Projektsicht erfüllt er damit nicht nur die Funktion einer Startoberfläche, sondern bildet die Verbindung zwischen zentral verwaltetem Spielbestand und tatsächlicher Nutzung auf dem Endgerät.

Die Beschreibung des Launchers lässt sich für den Bericht wie folgt zusammenfassen:
\begin{itemize}[leftmargin=1.5em]
  \item Der Launcher stellt die auf dem Gerät verfügbaren beziehungsweise zugewiesenen Android-Spiele in einer nutzbaren Oberfläche dar.
  \item Er dient als patientennahe Einstiegskomponente und reduziert den Aufwand, einzelne Spiele manuell zu suchen oder getrennt zu starten.
  \item Neue oder aktualisierte mobile Spiele werden nicht isoliert betrachtet, sondern in einen Plattformprozess aus Verwaltung, Zuweisung und Nutzung eingebettet.
  \item Für den Projektkontext ist der Launcher deshalb ein zentrales Bindeglied zwischen Web-/Backend-Verwaltung und realem Einsatz des Serious Games auf Android.
\end{itemize}

\subsection*{Einordnung für den Projektbericht}
Für diesen Bericht ist der Android-Launcher vor allem deshalb relevant, weil an ihm sichtbar wird, ob die Plattform ihr Ziel tatsächlich erreicht. Erst wenn ein Spiel nach Upload, Zuordnung und technischer Integration im Launcher erscheint und dort gestartet werden kann, wird aus einer Verwaltungsfunktion ein praktisch nutzbarer Therapie- oder Trainingsbaustein.

\subsection*{Optionaler Originalauszug aus der README}
\fbox{\parbox[c][6cm][c]{\textwidth}{\centering Platzhalter: Hier kann bei Bedarf der konkrete README-Auszug zum Android-Launcher eingefügt werden.}}

\end{document}
